\documentclass[b5paper,11pt,onecolumn,twoside,fontset=windows]{ctexart}
\pagestyle{empty}
\usepackage[total={130mm,203mm}, top=25.3mm, headheight=20pt, headsep=4truemm]{geometry}%调整版式
\usepackage{tabularx}
\usepackage{amssymb}
\usepackage{xeCJK}
\setmainfont{Times New Roman}

\newcommand{\vdistance}{\\[4.8ex]}%调整行间隔
\newcommand{\cndash}{\raise0.1em\hbox{--}}

\begin{document}

\begin{flushright}
\renewcommand{\baselinestretch}{1.3}
\zihao{5}
\heiti{逻辑学研究（双月刊）\\
第\ 19 卷，第\ 5 期，总第\ 90 期}\\
\heiti{2026 年第\ 5 期，10月\ 18 日出版}
\end{flushright}
\vspace{4ex}{\songti{\huge{\bfseries  目 \ \ \ \ \ \ \ \ 录}}}

\noindent\hrulefill \vspace{2.5ex}\\
\begin{tabularx}{\textwidth}
{l@{\hspace{-0.3em}}r@{\hspace{-4.2em}}r}	

%中间主体

{\large 论示例悖论及其解决方案}
\\{\large ——从扩展的X标准看}
&\textit{\large 陆王，王二}
& 1\\

\end{tabularx}
\end{document}
